\documentclass[10pt,a4paper]{amsart}
\usepackage[margin=19mm]{geometry}
\usepackage{amsmath,amssymb,mathtools}
\usepackage[scr=rsfs,cal=euler]{mathalfa}
\usepackage{xfrac}
\usepackage[unicode,hidelinks,bookmarks=false]{hyperref}
\newcommand{\ie}{i.e.\ }
\newcommand{\cf}{cf.\ }
\newcommand{\resp}{respectively\ }
\newcommand{\Subsection}{\subsection}
\providecommand{\nobreakdash}{\nobreak\hbox{-}}
\setlength{\emergencystretch}{2em}
\multlinegap0pt
\usepackage{bm}
\usepackage{bbm}
\usepackage{dsfont}
\usepackage{xspace}
\usepackage{cancel}
\usepackage{mathtools}
\usepackage{amsthm}
\makeatletter
\@ifundefined{lemma}{\newtheorem{lemma}{Lemma}}{}
\@ifundefined{theorem}{\newtheorem{theorem}{Theorem}}{}
\makeatother
\title[On Grundy indices for complete geometric graphs]{On Grundy indices for complete geometric graphs}
\author{Dolores Lara, Christian Rubio-Montiel, Francisco Zaragoza}
\date{Source: arXiv:2510.15155v1 (CC BY 4.0)}
\begin{document}
\maketitle

\noindent\emph{Source and licence.} On Grundy indices for complete geometric graphs. Version arXiv:2510.15155v1 (CC BY 4.0), distributed under the Creative Commons Attribution 4.0 International licence, as recorded in the arXiv licence metadata for this exact version and shown on the abstract page https://arxiv.org/abs/2510.15155v1. Full rights notice: licenses/arxiv-2510.15155-CC-BY-4.0.txt.\par\smallskip

\noindent\textit{Editorial scope.} Counted unit: the Theorem whose source label is \texttt{teo:remainder} in the article (source0.tex lines 501--504, source label \texttt{teo:remainder}), delivered together with the article's own complete original proof of it (source0.tex lines 505--526, one \texttt{proof} environment of 22 lines). No mathematics is written, repaired or re-derived: statement and proof are the article's own text line for line. Required context retained uncounted: lem:triangle-edge (lines 421--426); thm:partition (lines 431--433). Every definition, display and label of those lines is retained unchanged. Omitted from this package and not redistributed: 1--420, 427--430, 434--500, 527--550. Nothing inside a retained excerpt is omitted or altered: every retained body is delivered exactly as the article's lines are written, so no omission receipt of any kind accompanies this package. Citations: the retained excerpts carry no citation command at all, so no bibliography entry of the article is transcribed and no reference is redistributed. Reference adaptation: none was needed and none was made; every reference inside the delivered unit resolves inside it, so no unresolved reference or citation is produced. Printed slips kept literal: no printed word, symbol, hypothesis or number of the counted statement or of its proof was altered. Layout and class: the article's own class is \texttt{article} with packages including geometry, graphicx, amssymb, epstopdf, amsthm, amsmath, subcaption, comment, tikz, ifthen, soul, hyperref, xcolor; the wrapper is the lane's pinned \texttt{amsart} editorial wrapper at 10pt on A4, which restates the theorem-like environments the retained text needs and adds the article's own macro definitions that the retained text needs (none), with extra wrapper packages (bm, bbm, dsfont, xspace, cancel, mathtools). The delivered numbering is the unit's own. Assets: no retained span contains a figure environment, a graphics-inclusion command, a PDF-page inclusion, a table or a TikZ picture; no image file is copied into this package, the package declares no asset dependency and no image file is redistributed with it. Licence: version arXiv:2510.15155v1 of this article is distributed under the Creative Commons Attribution 4.0 International licence, as recorded in the arXiv licence metadata for this exact version and shown on the abstract page https://arxiv.org/abs/2510.15155v1. A full rights notice is delivered at licenses/arxiv-2510.15155-CC-BY-4.0.txt. Characters: every retained excerpt is pure ASCII, so the delivered engine, XeLaTeX with its default Latin Modern text font, reports no missing character.
\begin{lemma}\label{lem:triangle-edge}
Let $S$ be a set of points in general position in the plane. Let $T$ be any triangle with vertices in $S$, and $e$ be any segment that connects two points in $S$. Choose $T$ and $e$ so that they do not have vertices in common. $T$ contains at least one edge disjoint to $e$.
\begin{proof}
    The line expanded by the edge $e$ divides the plane into two half-planes; one half-plane contains exactly one vertex of $T$. The degree of this vertex is exactly two; therefore, there are at most two edges of $T$ that cross $e$.
\end{proof}
\end{lemma}

\begin{theorem}(H. Hanani)\label{thm:partition}
    Let $K_n$ be the complete graph with $n$ vertices. There exists a set $F \subset E(K_n)$ such that $E(K_n) \setminus F$ can be partitioned into triangles. More precisely, if $n\equiv 1, 3 \pmod 6$, then $F = \emptyset$, if $n\equiv 0, 2 \pmod 6$, then $F$ induces a perfect matching in $K_n$, if $n\equiv 4 \pmod 6$, then $F$ induces a spanning forest of $n/2 + 1$ edges in $K_n$ with all vertices having an odd degree (a tripole), and if $n\equiv 5 \pmod 6$, then $F$ induces a $4$-cycle. 
\end{theorem}

\begin{theorem}\label{teo:remainder} 
Let $K_n$ be the complete graph, under the non-crossing criteria the Grundy geometric index of $K_n$ has the following lower bound,
\[\frac{n^2}{6}+\Theta(n)\leq \Gamma'_g(K_n)\]
\end{theorem}

\begin{proof}

Let $S$ be a set of $n$ points in general position in the plane. Consider the set $F$ obtained by applying Theorem~\ref{thm:partition} to a complete graph $K_n$. Draw $F$ in $S$ so that there are no crossings between its elements. It is not hard to prove that such a drawing always exists, just order the points in $S$ from left to right, and let $S=\{s_1, \ldots, s_n\}$ be the ordered set. If $F$ is a perfect matching draw the edges $(s_1,s_2), (s_2,s_3), \ldots, (s_{n-1}, s_n)$; if $F$ is a tripole, choose the first three points and draw the $K_{1,3}$ graph, and then draw the perfect matching $(s_4,s_5), (s_6,s_7), \ldots, (s_{n-1}, s_n)$; finally, if $F$ is a $4$-cycle, just draw the edges $(s_1,s_2), (s_2,s_3), (s_3,s_4), (s_4,s_1)$.

Without loss of generality, let $\mathsf{K}_n$ be the complete graph drawn in the plane as described above. Theorem~\ref{thm:partition} guarantees that the set of edges of $\mathsf{K}_n \setminus F$ can be partitioned into $k$ triangles. Color each of these triangles with a different color, using colors $\{1, 2, \ldots, k\}$. 
This partial coloring is pseudo-Grundy, since Lemma~\ref{lem:triangle-edge} implies that for every two triangles in the partition there are always two parallel edges, one from each triangle; however, the coloring is not proper. Now we must color $F$.

Color each edge in $F$ with a different color, using colors $\{k+1, \ldots, k+|F|\}$. By construction, every pair of edges in $F$ is parallel or shares exactly one vertex; therefore, $F$ induces a clique. This implies that the coloring of $F$ is Grundy. Now consider an edge $e$ in $F$ colored with color $j$, and let $T$ be a triangle in the partition, colored with color $i$, $i < j$. There are two possible situations; either $e$ and $T$ do not share any vertex or $e$ is incident to one vertex of $T$. In the first case, Lemma~\ref{lem:triangle-edge} guarantees that there is at least one edge in $T$ parallel to $e$, and its chromatic classes are adjacent. In the latter case, the condition is met by definition. Therefore, the coloring is pseudo-Grundy.


The number of colors used is 

\begin{equation*}
 k + |F|  = 
  \begin{cases}
    \frac{(n-1)n}{6} & \text{if } n\equiv 1, 3 \pmod 6\\
    \frac{(n+1)n}{6} & \text{if } n\equiv 0, 2 \pmod 6\\
     \frac{(n-1)n+16}{6}  & \text{if } n\equiv 4 \pmod 6\\
    \frac{(n+1)n+4}{6} & \text{if } n\equiv 5 \pmod 6
  \end{cases}
\end{equation*}
\end{proof}

\end{document}
